\documentclass[10pt,a4paper]{amsart}
\usepackage[margin=19mm]{geometry}
\usepackage{amsmath,amssymb,mathtools}
\usepackage[scr=rsfs,cal=euler]{mathalfa}
\usepackage{xfrac}
\usepackage[unicode,hidelinks,bookmarks=false]{hyperref}
\newcommand{\ie}{i.e.\ }
\newcommand{\cf}{cf.\ }
\newcommand{\resp}{respectively\ }
\newcommand{\Subsection}{\subsection}
\providecommand{\nobreakdash}{\nobreak\hbox{-}}
\setlength{\emergencystretch}{2em}
\multlinegap0pt
\usepackage{bm}
\usepackage{bbm}
\usepackage{dsfont}
\usepackage{xspace}
\usepackage{cancel}
\usepackage{mathtools}
\usepackage{amsthm}
\makeatletter
\@ifundefined{theorem}{\newtheorem{theorem}{Theorem}}{}
\makeatother
\title[Composition operators and Rational Inner Functions on the bi]{Composition operators and Rational Inner Functions on the bidisc: A geometric approach}
\author{Athanasios Beslikas}
\date{Source: arXiv:2604.20330v1 (CC BY 4.0)}
\begin{document}
\maketitle

\noindent\emph{Source and licence.} Composition operators and Rational Inner Functions on the bidisc: A geometric approach. Version arXiv:2604.20330v1 (CC BY 4.0), distributed under the Creative Commons Attribution 4.0 International licence, as recorded in the arXiv licence metadata for this exact version and shown on the abstract page https://arxiv.org/abs/2604.20330v1. Full rights notice: licenses/arxiv-2604.20330-CC-BY-4.0.txt.\par\smallskip

\noindent\textit{Editorial scope.} Counted unit: the Theorem whose source label is \texttt{None} in the article (source0.tex lines 457--460, source label \texttt{None}), delivered together with the article's own complete original proof of it (source0.tex lines 461--486, one \texttt{proof} environment of 26 lines). No mathematics is written, repaired or re-derived: statement and proof are the article's own text line for line. Required context retained uncounted: none. Every definition, display and label of those lines is retained unchanged. Omitted from this package and not redistributed: 1--456, 487--998. Nothing inside a retained excerpt is omitted or altered: every retained body is delivered exactly as the article's lines are written, so no omission receipt of any kind accompanies this package. Citations: the retained excerpts carry no citation command at all, so no bibliography entry of the article is transcribed and no reference is redistributed. Reference adaptation: none was needed and none was made; every reference inside the delivered unit resolves inside it, so no unresolved reference or citation is produced. Printed slips kept literal: no printed word, symbol, hypothesis or number of the counted statement or of its proof was altered. Layout and class: the article's own class is \texttt{amsart} with packages including geometry, fullpage, framed, xcolor, tikz, amsmath, graphicx, amsfonts, amssymb, hyperref, etoolbox, caption, titletoc; the wrapper is the lane's pinned \texttt{amsart} editorial wrapper at 10pt on A4, which restates the theorem-like environments the retained text needs and adds the article's own macro definitions that the retained text needs (none), with extra wrapper packages (bm, bbm, dsfont, xspace, cancel, mathtools). The delivered numbering is the unit's own. Assets: no retained span contains a figure environment, a graphics-inclusion command, a PDF-page inclusion, a table or a TikZ picture; no image file is copied into this package, the package declares no asset dependency and no image file is redistributed with it. Licence: version arXiv:2604.20330v1 of this article is distributed under the Creative Commons Attribution 4.0 International licence, as recorded in the arXiv licence metadata for this exact version and shown on the abstract page https://arxiv.org/abs/2604.20330v1. A full rights notice is delivered at licenses/arxiv-2604.20330-CC-BY-4.0.txt. Characters: every retained excerpt is pure ASCII, so the delivered engine, XeLaTeX with its default Latin Modern text font, reports no missing character.
\begin{theorem}\textbf{(Smooth tangential intersection)} Let $\Phi=(\phi_1,\phi_2)$ where $\phi_1,\phi_2$ RIFs and assume that there exist two analytic curves $(\zeta,g_{j_1}^{\phi_1,\alpha_1}(\zeta))\in\mathcal{C}_{{\alpha}_1}(\phi_1)$ and $(\zeta,g^{\alpha_2,\phi_2}_{j_2}(\zeta))\in\mathcal{C}_{\alpha_2}(\phi_2),$ for some $\alpha_1,\alpha_2\in\mathbb{T}$ such that they intersect tangentially at a point $\zeta_0\in\mathbb{T}^2$ away from the (possible) singularities, i.e. there exists an exponent $M>1$ such that 
$$|g^{\alpha_1,\phi_1}_{j_1}(\zeta)-g_{j_2}^{\alpha_2,\phi_2}(\zeta)|\approx |\zeta-\zeta_0|^M$$
Then, the composition operator $C_{\Phi}:A^2_{\beta}(\mathbb{D}^2)\to A^2_{\beta}(\mathbb{D}^2)$ is not bounded. 
\end{theorem}

\begin{proof} The strategy for the proof of our previous result will be followed after some necessary changes. We choose to include this proof separately to emphasize here that the two curves do not share an arc; they meet only on a point of tangential intersection. Moreover, the tangential intersection is away from the singularity, therefore, we can apply local arguments on the partials and take $\partial_{z_i}\phi_{\ell}(\zeta,g_{j_\ell}^{\alpha_{\ell},\phi_{\ell}}(\zeta))=c_{i\ell}\neq0,$ for $\ell=1,2.$ Therefore, similar estimates to our previous result hold. More specifically

By applying the same ideas, we obtain two local estimates

$$|\phi_1(r\zeta,z_2)-\alpha_1|\lesssim (1-r)+|z_2-g^{\alpha_1,\phi_1}_{j_1}(\zeta)|,$$
and
\begin{align}
|\phi_2(r\zeta,z_2)-\alpha_2|&\lesssim (1-r)+|z_2-g^{\phi_2,\alpha_2}_{j_2}(\zeta)| \notag&&\\
&\lesssim(1-r)+|g_{j_1}^{\phi_1,\alpha_1}(\zeta)-g_{j_2}^{\phi_2,\alpha_2}(\zeta)|+|z_2-g_{j_1}^{\phi_1,\alpha_1}(\zeta)|\notag&&\\
&\approx (1-r)+|\zeta-\zeta_0|^M+|z_2-g^{\phi_1,\alpha_1}_{j_1}(\zeta)|
.\end{align}
In a similar manner to the previous proof we need to consider the set

$$A_{\delta}:=\Bigg\{(r\zeta,z_2)\in\mathbb{D}^2:\zeta\in I,\quad|\zeta-\zeta_0|^M<C_1\delta,\quad 1-r<C_2\delta, \quad|z_2-g_{j_1}^{\phi_1,\alpha_1}(\zeta)|<\delta\Bigg\},$$

and observe $A_{\delta}\subset \Phi^{-1}(S((\alpha_1,\alpha_2),\delta)).$
We will now estimate the volume of $A_{\delta}.$

\begin{align}
V_{\beta}(\Phi^{-1}(S((\alpha_1,\alpha_2),\delta)))>V_{\delta}(A_{\beta})=&\int_{\{(r\zeta,z_1)\in\mathbb{D}^2:\zeta\in I,|\zeta-\zeta_0|^M<C_1\delta,1-r<C_2\delta, |z_2-g^{\alpha_1,\phi_1}_{j_1}(\zeta)|<\delta\}}dV_{\beta}(z_1,z_2)\notag&&\\
=&\left(\int_{\{|\zeta-\zeta_0|^M<C_1\delta,1-r<C_2\delta\}}dA_{\beta}(z_1)\right)\left(\int_{|z_2-g^{\alpha_1,\phi_1}_{j_1}(\zeta)|<\delta}dA_{\beta}(z_2)\right)\notag&&\\
\approx &\delta^{1/M}\delta^{\beta+1}\delta^{\beta+2}&&\\
\approx &\delta^{2\beta+3+\frac{1}{M}} \notag
\end{align}
Since $M>1,$ this suffices to show non-boundedness of the composition operator.
\end{proof}

\end{document}
